% ==================== algorithms ====================
\section{Алгоритмы перехода}
\label{sec:algorithms}

Теперь применим условия главы~\ref{sec:optimality} к расчёту манёвров.
Для каждого типа решения будем выполнять одинаковые действия:
устанавливать геометрию годографа, находить импульсы из граничных
уравнений и проверять, что достигнута нижняя оценка затрат.
Такой порядок соответствует геометрическому изложению Баранова
\cite[гл.~2]{baranov-ru} и классификации Эдельбаума~\cite{edelbaum}.

Исходные данные --- $\Delta a,\Delta e_x,\Delta e_y,\Delta i$ из
раздела~\ref{sec:five-boundaries}. Напомним:
\begin{equation}
 \Delta a=\frac{a_f-a_i}{r_0},\qquad
 \Delta e=\sqrt{(\Delta e_x)^2+(\Delta e_y)^2},\qquad \Delta i\geq0.
 \label{eq:transfer-data}
\end{equation}
Здесь $\Delta e$ --- длина \emph{разности векторов} эксцентриситета,
а не разность чисел $e_f-e_i$. Все компоненты импульсов ниже безразмерны.
Для получения м/с их умножают на $V_0$, выраженную в м/с.
Число $J$ есть сумма модулей, поэтому размерная характеристическая
скорость равна $\Delta V_\Sigma=V_0J$.

\subsection{Компланарный случай}
\label{sec:planar}

\subsubsection{Какие изменения создаёт трансверсальный импульс}

При $\Delta i=0$ нормальные импульсы не нужны: их удаление
не меняет первые три граничных условия и уменьшает либо сохраняет $J$.
Остаются радиальные и трансверсальные компоненты.
Из анализа плоского годографа в разделе~\ref{sec:hodograph}
получены схемы с трансверсальными импульсами.
Сначала построим такие схемы, затем проверим их оптимальность
среди всех импульсов, в том числе имеющих радиальные компоненты.

При $\Delta\widehat V_{rj}=0$ условия перехода имеют вид
\begin{equation}
 \begin{aligned}
 2\sum_j\Delta\widehat V_{tj}&=\Delta a,\\
 2\sum_j\Delta\widehat V_{tj}\cos\varphi_j&=\Delta e_x,\\
 2\sum_j\Delta\widehat V_{tj}\sin\varphi_j&=\Delta e_y.
 \end{aligned}
 \label{eq:planar-boundaries}
\end{equation}
Один трансверсальный импульс изменяет большую полуось на
$2\Delta\widehat V_t$ в принятой нормировке и добавляет к вектору
эксцентриситета вектор
\[
 2\Delta\widehat V_t(\cos\varphi,\sin\varphi).
\]
Это позволяет изображать переход на плоскости $(e_x,e_y)$.
Точка на ней задаёт вектор эксцентриситета орбиты, а отрезок между
двумя точками --- его изменение. Большая полуось учитывается
первой строкой~\eqref{eq:planar-boundaries}. Длина отрезка равна
$2|\Delta\widehat V_t|$. Для разгонного импульса направление отрезка
совпадает с углом $\varphi$; для тормозного --- противоположно ему.
Эта диаграмма не изображает траекторию аппарата в физическом пространстве.

\subsubsection{Два импульса через полвитка}

Если $\Delta e>0$, определим направление требуемого изменения
эксцентриситета:
\begin{equation}
 \cos\varphi_e=\frac{\Delta e_x}{\Delta e},\qquad
 \sin\varphi_e=\frac{\Delta e_y}{\Delta e}.
 \label{eq:phi-e}
\end{equation}
Численно это $\varphi_e=\operatorname{atan2}(\Delta e_y,\Delta e_x)$.
Функция $\operatorname{atan2}$ определяет угол с учётом знаков обеих
компонент; обычный арктангенс их отношения не различает противоположные
направления.

Выберем места импульсов
$\varphi_1=\varphi_e$, $\varphi_2=\varphi_e+\pi$.
У второго импульса синус и косинус меняют знак. Поэтому из трёх
уравнений~\eqref{eq:planar-boundaries} остаются два независимых:
\[
 2(\Delta\widehat V_{t1}+\Delta\widehat V_{t2})=\Delta a,
 \qquad
 2(\Delta\widehat V_{t1}-\Delta\widehat V_{t2})=\Delta e.
\]
Сложение и вычитание дают
\begin{equation}
 \boxed{\begin{aligned}
 \varphi_1&=\varphi_e,&\varphi_2&=\varphi_e+\pi,\\
 \Delta\widehat V_{t1}&=\frac{\Delta a+\Delta e}{4},&
 \Delta\widehat V_{t2}&=\frac{\Delta a-\Delta e}{4},\\
 \Delta\widehat V_{r1}&=0,&\Delta\widehat V_{r2}&=0.
 \end{aligned}}
 \label{eq:planar-impulses}
\end{equation}
Это \emph{компоненты}, поэтому они могут быть отрицательными.
Стоимость найденного перехода равна
\begin{equation}
 \boxed{\ J=\frac{|\Delta a+\Delta e|+|\Delta a-\Delta e|}{4}
       =\frac12\max\{|\Delta a|,\Delta e\}.\ }
 \label{eq:planar-cost}
\end{equation}
Последнее равенство проверяется отдельно при
$|\Delta a|\geq\Delta e$ и $|\Delta a|\leq\Delta e$.
Направление $\varphi_e$ задаёт линия апсид \emph{относительной орбиты}:
оно определяется разностью векторов эксцентриситета и не обязано
совпадать с линией апсид какой-либо из двух исходных орбит.

\subsubsection{Почему меньших затрат быть не может}

Первое граничное условие при любых импульсах даёт
\[
 |\Delta a|=2\left|\sum_j\Delta\widehat V_{tj}\right|
      \leq2\sum_j|\Delta\widehat V_{tj}|\leq2J.
\]
Рассмотрим также изменение вектора эксцентриситета от одного импульса.
Сумма квадратов его двух компонент равна
\[
 (\Delta\widehat V_{rj}\sin\varphi_j+2\Delta\widehat V_{tj}\cos\varphi_j)^2
 +(-\Delta\widehat V_{rj}\cos\varphi_j+2\Delta\widehat V_{tj}\sin\varphi_j)^2
 =(\Delta\widehat V_{rj})^2+4(\Delta\widehat V_{tj})^2.
\]
Смешанные слагаемые сокращаются, а $\sin^2\varphi_j+\cos^2\varphi_j=1$.
Длина этого изменения не больше $2\Delta\widehat V_j$.
Длина суммы векторов не превосходит суммы их длин, следовательно,
$\Delta e\leq2J$. Вместе получаем
\begin{equation}
 J\geq\frac12\max\{|\Delta a|,\Delta e\}.
 \label{eq:planar-bound}
\end{equation}
Схема~\eqref{eq:planar-impulses} достигает этой оценки.
Значит, она оптимальна при любом разрешённом числе импульсов.
Радиальные компоненты при доказательстве не запрещались.

\subsubsection{Пересечение орбит и вид годографа}

В первом порядке разность радиусов заданных орбит в одной фазе равна
\begin{equation}
 \frac{r_f(\varphi)-r_i(\varphi)}{r_0}
 =\Delta a-\Delta e_x\cos\varphi-\Delta e_y\sin\varphi
 =\Delta a-\Delta e\cos(\varphi-\varphi_e).
 \label{eq:planar-intersection}
\end{equation}
Отсюда следуют три случая. Пересечение и касание здесь понимаются
в геометрии первого приближения.

\paragraph{Пересекающиеся орбиты: $|\Delta a|<\Delta e$.}
Разность радиусов обращается в нуль при двух фазах.
Первый импульс в~\eqref{eq:planar-impulses} разгонный,
второй тормозной. Их вклады в вектор эксцентриситета сонаправлены,
поскольку одновременно меняются знак импульса и направление радиуса.
Стоимость равна $J=\Delta e/2$.

Этой схеме отвечает центральный эллипс
\begin{equation}
 s_r=\tfrac12\sin(\varphi-\varphi_e),\qquad
 s_t=\cos(\varphi-\varphi_e),\qquad s_z=0.
 \label{eq:planar-central-primer}
\end{equation}
Действительно, $p^2=1-\tfrac34\sin^2(\varphi-\varphi_e)\leq1$.
В двух местах импульсов $\boldsymbol s=(0,1,0)$ и $(0,-1,0)$.
Они совпадают с направлениями найденных импульсов.
Геометрическая и прямая проверки оптимальности дают одинаковый результат.
Импульсы, вообще говоря, прикладываются не в точках пересечения орбит.
Возможность перейти одним импульсом в точке пересечения ещё не означает
минимума затрат.

\paragraph{Непересекающиеся орбиты: $|\Delta a|>\Delta e$.}
Разность радиусов не меняет знак. Оба импульса
в~\eqref{eq:planar-impulses} имеют знак $\Delta a$:
при увеличении большой полуоси оба разгонные, при уменьшении ---
оба тормозные. Стоимость $J=|\Delta a|/2$.
Годограф вырождается в точку
\[
 \boldsymbol s=(0,\operatorname{sgn}\Delta a,0).
\]
Условие единичного модуля выполнено при любой фазе.
Поэтому противоположные места --- удобный выбор, но не единственный.

\paragraph{Касание орбит: $|\Delta a|=\Delta e>0$.}
Один импульс в~\eqref{eq:planar-impulses} исчезает.
Остаётся трансверсальный импульс в точке касания:
$\varphi=\varphi_e$ при $\Delta a>0$ и
$\varphi=\varphi_e+\pi$ при $\Delta a<0$.
Его компонента удовлетворяет
$\Delta\widehat V_t=\Delta a/2$.
Смещённый эллипс с единственным касанием на рис.~\ref{fig:planar-hodographs}
соответствует именно такой геометрии перехода.

\subsubsection{Графический метод на плоскости эксцентриситетов}

Обозначим на этой плоскости начальную, конечную и промежуточную
орбиты точками $A$, $K$ и $M$. Для двух трансверсальных импульсов
\[
 AM=2|\Delta\widehat V_{t1}|,\qquad
 MK=2|\Delta\widehat V_{t2}|,\qquad AK=\Delta e.
\]
Если оба импульса разгонные, их сумма фиксирована условием на
большую полуось, поэтому $AM+MK=\Delta a$.
При двух тормозных импульсах так же получается
$AM+MK=|\Delta a|$. Значит, при $|\Delta a|>\Delta e$
точка $M$ лежит на эллипсе с фокусами $A,K$.
Каждая точка этого эллипса задаёт оптимальный двухимпульсный переход
той же стоимости $|\Delta a|/2$.

Если $\Delta e>|\Delta a|$ и первый импульс разгонный, а второй
тормозной, условие на большую полуось имеет вид $AM-MK=\Delta a$.
При $\Delta a\ne0$ это одна из ветвей гиперболы; при $\Delta a=0$
получается серединный перпендикуляр к $AK$. Стоимость равна
$(AM+MK)/2$ и меняется
при перемещении $M$ по ветви. Наименьшая сумма достигается,
когда $M$ лежит на отрезке $AK$: тогда $AM+MK=AK=\Delta e$.
Совместно с разностью длин это снова даёт~\eqref{eq:planar-impulses}.

\begin{figure}[htbp]
\centering
\begin{tikzpicture}[x=2cm,y=2cm,>=Latex]
 \begin{scope}
  \draw[->,gray] (-1.6,-1.15)--(1.65,-1.15) node[right] {$e_x$};
  \draw[->,gray] (-1.6,-1.15)--(-1.6,1.35) node[above] {$e_y$};
  \draw[accent,thick] (0,0) ellipse[x radius=1.25,y radius=.9605];
  \coordinate (A) at (-.8,0); \coordinate (K) at (.8,0); \coordinate (M) at (0,.9605);
  \draw[->,very thick,accent] (A)--(M); \draw[->,very thick,accent] (M)--(K);
  \draw[dashed] (A)--(K);
  \fill (A) circle(1.5pt) node[below] {$A$};
  \fill (K) circle(1.5pt) node[below] {$K$};
  \fill (M) circle(1.5pt) node[above] {$M$};
  \node at (0,-1.5) {$AM+MK=|\Delta a|$};
 \end{scope}
 \begin{scope}[shift={(4.5,0)}]
  \draw[->,gray] (-1.6,-1.15)--(1.65,-1.15) node[right] {$e_x$};
  \draw[->,gray] (-1.6,-1.15)--(-1.6,1.35) node[above] {$e_y$};
  \draw[accent,thick,domain=-1:1,samples=60] plot ({.35*sqrt(1+\x*\x/.6875)},\x);
  \coordinate (A) at (-.9,0); \coordinate (K) at (.9,0); \coordinate (M) at (.35,0);
  \coordinate (P) at ({.35*sqrt(1+.81/.6875)},.9);
  \draw[dashed,gray] (A)--(P)--(K);
  \draw[->,very thick,accent] (A)--(M);\draw[->,very thick,accent] (M)--(K);
  \fill (A) circle(1.5pt) node[below] {$A$};
  \fill (M) circle(1.5pt) node[below] {$M$};
  \fill (K) circle(1.5pt) node[below] {$K$};
  \fill (P) circle(1.5pt) node[above] {$M'$};
  \node at (0,-1.5) {$AM-MK=\Delta a$};
 \end{scope}
\end{tikzpicture}
\caption{Метод Баранова на плоскости эксцентриситетов, показанный
для $\Delta e_y=0$, $\Delta a>0$. Слева все положения $M$ на эллипсе
равноценны. Справа оптимально положение $M$ на отрезке $AK$;
ломаная через $M'$ длиннее. Отрезки изображают изменения вектора
эксцентриситета, а не перемещения аппарата.}
\label{fig:eccentricity-geometry}
\end{figure}

\paragraph{Если место первого импульса выбрано заранее.}
Для непересекающихся орбит можно задать любую фазу $\varphi_1$.
Тогда из~\eqref{eq:planar-boundaries}
\[
 2\Delta\widehat V_{t2}\cos\varphi_2
      =\Delta e_x-2\Delta\widehat V_{t1}\cos\varphi_1,\qquad
 2\Delta\widehat V_{t2}\sin\varphi_2
      =\Delta e_y-2\Delta\widehat V_{t1}\sin\varphi_1.
\]
Возведём эти равенства в квадрат, сложим и подставим
$\Delta\widehat V_{t2}=\Delta a/2-\Delta\widehat V_{t1}$.
Квадраты неизвестной компоненты сокращаются. Получаем формулы Баранова:
\begin{equation}
\begin{aligned}
 \Delta\widehat V_{t1}&=
 \frac{\Delta e^2-(\Delta a)^2}
 {4(\Delta e_x\cos\varphi_1+\Delta e_y\sin\varphi_1-\Delta a)},\\
 \Delta\widehat V_{t2}&=\frac{\Delta a}{2}-\Delta\widehat V_{t1},\\
 \cos\varphi_2&=\frac{\Delta e_x-2\Delta\widehat V_{t1}\cos\varphi_1}
                         {2\Delta\widehat V_{t2}},\qquad
 \sin\varphi_2=\frac{\Delta e_y-2\Delta\widehat V_{t1}\sin\varphi_1}
                         {2\Delta\widehat V_{t2}}.
\end{aligned}
\label{eq:planar-free-first}
\end{equation}
При $|\Delta a|>\Delta e$ знаменатели не равны нулю,
а обе компоненты имеют знак $\Delta a$. По последним двум строкам
восстанавливают угол с учётом квадранта, включая знак
$\Delta\widehat V_{t2}$. При касании орбит пользуются одноимпульсной
формулой, а не делят на исчезающий второй импульс.

\subsubsection{Пример и простые предельные случаи}

Пусть $\Delta a=0{,}001$, $\Delta e_x=0{,}002$,
$\Delta e_y=0{,}001$. Тогда
\[
 \Delta e=0{,}00223607,\qquad \varphi_e=26{,}5651^\circ,
 \qquad\Delta\widehat V_{t1}=0{,}000809017,
 \quad\Delta\widehat V_{t2}=-0{,}000309017.
\]
Второй импульс прикладывается при $206{,}5651^\circ$.
При $r_0=7000$ км и $\mu=398600{,}4418$ км$^3$/с$^2$
имеем $V_0=7546{,}053$ м/с. Размерные компоненты равны
$6{,}105$ и $-2{,}332$ м/с, а их сумма модулей --- $8{,}437$ м/с.
Условие на большую полуось проверяется сложением компонентов,
условие на эксцентриситет --- их разностью:
\[
 2(\Delta\widehat V_{t1}+\Delta\widehat V_{t2})=0{,}001,
 \qquad 2(\Delta\widehat V_{t1}-\Delta\widehat V_{t2})=\Delta e.
\]

При $\Delta e=0$ угол $\varphi_e$ не определён и вычислять его
не нужно. Достаточно любых двух противоположных фаз и импульсов
$\Delta\widehat V_{t1}=\Delta\widehat V_{t2}=\Delta a/4$.
Это линейный аналог перехода между близкими круговыми орбитами.
При $\Delta a=0$, $\Delta e>0$ компоненты равны $\Delta e/4$
и $-\Delta e/4$. При $\Delta a=\Delta e=0$ манёвры не требуются.

% ==================== spatial ====================
\subsection{Некомпланарный случай}
\label{sec:spatial}

\subsubsection{Выбор осей и геометрия решений}

Теперь $\Delta i>0$. Как в разделе~\ref{sec:five-boundaries},
ось $X$ направлена вдоль линии пересечения плоскостей к восходящему
узлу конечной орбиты относительно начальной. Начальная плоскость
служит опорной. В этих осях заданы $\Delta e_x,\Delta e_y$, а условия
на нормальные импульсы равны
\begin{equation}
 \sum_j\Delta\widehat V_{zj}\sin\varphi_j=0,\qquad
 \sum_j\Delta\widehat V_{zj}\cos\varphi_j=\Delta i.
 \label{eq:plane-change-boundaries}
\end{equation}
Если исходные плоскости заданы наклонениями $i_i,i_f$ и долготами
узлов $\Omega_i,\Omega_f$, то
\[
 \cos\Delta i=\cos i_i\cos i_f+
       \sin i_i\sin i_f\cos(\Omega_f-\Omega_i).
\]
В общем случае $\Delta i$ не равно $|i_f-i_i|$.
Узловой базис можно построить по единичным нормалям
$\boldsymbol k_i,\boldsymbol k_f$:
\[
 \boldsymbol e_X=\frac{\boldsymbol k_i\times\boldsymbol k_f}
                         {\|\boldsymbol k_i\times\boldsymbol k_f\|},
 \qquad\boldsymbol e_Z=\boldsymbol k_i,\qquad
 \boldsymbol e_Y=\boldsymbol e_Z\times\boldsymbol e_X.
\]
Проекции $\boldsymbol e_f-\boldsymbol e_i$ на $\boldsymbol e_X,\boldsymbol e_Y$
дают нужные $\Delta e_x,\Delta e_y$.

На рис.~\ref{fig:spatial-hodographs} показаны три геометрические
конфигурации, рассмотренные в главе~2.
Теперь для каждой из них получим направления и величины импульсов.

\subsubsection{Узловое решение: сначала направления, затем величины}
\label{sec:NC}

У центрального эллипса противоположные касания соответствуют фазам,
различающимся на $\pi$, и противоположным локальным направлениям.
Пусть первое направление равно
$\boldsymbol d=(d_r,d_t,d_z)$, $\|\boldsymbol d\|=1$.
Тогда импульсы имеют вид
$\Delta\widehat{\boldsymbol V}_1=\Delta\widehat V_1\boldsymbol d$,
$\Delta\widehat{\boldsymbol V}_2=-\Delta\widehat V_2\boldsymbol d$.
Их нормальный вклад в~\eqref{eq:plane-change-boundaries} пропорционален
$Jd_z(\sin\varphi_1,\cos\varphi_1)$. При $\Delta i>0$ первая
компонента должна быть нулевой, значит $\sin\varphi_1=0$.
Можно нумеровать места как $\varphi_1=0$, $\varphi_2=\pi$.
Так из геометрии годографа получается именно \emph{узловое} решение.

Подставим эти углы и направления в пять граничных уравнений:
\begin{equation}
\begin{aligned}
 \Delta a&=2d_t(\Delta\widehat V_1-\Delta\widehat V_2),\\
 \Delta e_x&=2d_t(\Delta\widehat V_1+\Delta\widehat V_2)=2Jd_t,\\
 \Delta e_y&=-d_r(\Delta\widehat V_1+\Delta\widehat V_2)=-Jd_r,\\
 \Delta i&=d_z(\Delta\widehat V_1+\Delta\widehat V_2)=Jd_z.
\end{aligned}
\label{eq:nc-balances}
\end{equation}
Условие с синусом выполнено автоматически. Из трёх последних строк
находим направляющие косинусы. Их квадраты должны давать единицу:
\begin{equation}
 \boxed{\ J_{\rm NC}=\sqrt{\frac{(\Delta e_x)^2}{4}+(\Delta e_y)^2+(\Delta i)^2},
 \qquad \boldsymbol d=\frac{1}{J_{\rm NC}}
       \begin{pmatrix}-\Delta e_y\\\Delta e_x/2\\\Delta i\end{pmatrix}.\ }
 \label{eq:nc-cost}
\end{equation}
Величины импульсов теперь находятся по их сумме и разности:
\begin{equation}
 \Delta\widehat V_1=\frac{J_{\rm NC}}2\left(1+\frac{\Delta a}{\Delta e_x}\right),
 \qquad
 \Delta\widehat V_2=\frac{J_{\rm NC}}2\left(1-\frac{\Delta a}{\Delta e_x}\right).
 \label{eq:nc-magnitudes}
\end{equation}
Пока предполагается $\Delta e_x\ne0$.
Модули не могут быть отрицательными, поэтому необходимо
$|\Delta a|\leq|\Delta e_x|$. Это первое ограничение области NC.

\paragraph{Проверка годографа.}
У центрального эллипса $\lambda_1=0$. Требование
$\boldsymbol s(0)=\boldsymbol d$ определяет три множителя:
$\lambda_3=d_t/2$, $\lambda_2=-d_r$, $\lambda_5=d_z$.
Четвёртый находится из касания в узле:
\[
 \boldsymbol s(0)\cdot\boldsymbol s'(0)
       =3\lambda_3\lambda_2+\lambda_5\lambda_4=0,
 \qquad \lambda_4=-\frac{3\lambda_3\lambda_2}{\lambda_5}.
\]
Здесь $\lambda_5\ne0$, поскольку $\Delta i>0$.
После подстановки получаем
\begin{equation}
 \|\boldsymbol s(\varphi)\|^2
 =1-\frac{\bigl[(\Delta i)^2-3(\Delta e_y)^2\bigr]
          \bigl[1+3(\Delta e_x)^2/(16(\Delta i)^2)\bigr]}
             {J_{\rm NC}^2}\sin^2\varphi.
 \label{eq:nc-norm}
\end{equation}

\begin{explanation}{Как получить и прочитать~\eqref{eq:nc-norm}}
При $\lambda_1=0$ можно сгруппировать синусы и косинусы:
\[
 \boldsymbol s(\varphi)=\boldsymbol d\cos\varphi+
             (\lambda_3,2\lambda_2,\lambda_4)\sin\varphi.
\]
Два постоянных вектора ортогональны благодаря условию касания.
Поэтому квадрат модуля равен
$\cos^2\varphi+(\lambda_3^2+4\lambda_2^2+\lambda_4^2)\sin^2\varphi$.
Подстановка найденных множителей и выделение $1$ дают
\eqref{eq:nc-norm}. Если коэффициент вычитаемого слагаемого
неотрицателен, эллипс целиком лежит внутри единичной сферы.
\end{explanation}

Таким образом, узловое решение оптимально при
\begin{equation}
 \boxed{\quad |\Delta a|\leq|\Delta e_x|,\qquad
                   \Delta i\geq\sqrt3\,|\Delta e_y|.\quad}
 \label{eq:nc-region}
\end{equation}
Оба условия имеют понятное происхождение: первое обеспечивает
неотрицательные модули, второе --- допустимый годограф.
Все граничные условия уже выполнены, поэтому применима проверка
из раздела~\ref{sec:optimality-check-full} при $\lambda_6=0$.

Если $\Delta e_x=0$, условие области требует $\Delta a=0$.
Делить на $\Delta e_x$ в~\eqref{eq:nc-magnitudes} нельзя.
Достаточен один импульс $(-\Delta e_y,0,\Delta i)$ при $\varphi=0$.
Его можно также разделить между двумя узлами с противоположными
локальными направлениями без изменения суммы модулей.

\paragraph{Численный пример.}
Пусть $\Delta a=0{,}001$, $\Delta e_x=0{,}002$,
$\Delta e_y=0{,}0005$, $\Delta i=0{,}002$.
Оба неравенства~\eqref{eq:nc-region} выполнены. Получаем
\[
 \begin{gathered}
 J=0{,}00229129,\\
 \Delta\widehat{\boldsymbol V}_1=(-0{,}000375,\ 0{,}000750,\ 0{,}001500),\\
 \Delta\widehat{\boldsymbol V}_2=(0{,}000125,\ -0{,}000250,\ -0{,}000500).
 \end{gathered}
\]
Компоненты каждого вектора перечислены в порядке $(r,t,z)$.
Например, $2(\Delta\widehat V_{t1}-\Delta\widehat V_{t2})=0{,}002$,
$-\Delta\widehat V_{r1}+\Delta\widehat V_{r2}=0{,}0005$,
$\Delta\widehat V_{z1}-\Delta\widehat V_{z2}=0{,}002$.
При принятом выше $V_0$ размерная стоимость равна $17{,}290$ м/с.

\subsubsection{Смещённый эллипс: симметрия двух импульсов}
\label{sec:NDC}

Рассмотрим невырожденное решение NDC при $\Delta e_y\ne0$.
Годограф имеет две симметричные точки касания, как на
рис.~\ref{fig:spatial-hodographs}. Его удобно записать в виде
\begin{equation}
 s_r=R\sin\psi,\qquad s_t=D+2R\cos\psi,\qquad
 s_z=F\sin\psi,\qquad \psi=\varphi-\alpha.
 \label{eq:ndc-hodograph}
\end{equation}
Здесь $R>0$; $D$ --- смещение центра вдоль $s_t$.
В отличие от узлового случая, два касания не противоположны.
Пусть им отвечают $\psi=\beta$ и $\psi=-\beta$, где $0<\beta<\pi$.
Тогда
\begin{equation}
 \varphi_+=\alpha+\beta,\qquad\varphi_-=\alpha-\beta,
 \qquad
 \boldsymbol d_+=(d_r,d_t,d_z),\quad
 \boldsymbol d_-=(-d_r,d_t,-d_z).
 \label{eq:ndc-symmetry}
\end{equation}
Знаки $+$ и $-$ обозначают стороны годографа, а не порядок манёвров
во времени. На опорной орбите две выбранные фазы симметричны
относительно направления $\alpha$, отсчитанного от узловой оси.
Угол $\beta$ задаёт их смещение от этого направления.
Это вспомогательные углы расчёта, не новые элементы орбиты.

Сначала выразим условие касания через направления.
Дифференцируя квадрат модуля~\eqref{eq:ndc-hodograph} при $\psi=\beta$,
получаем
\[
 (R^2+F^2)\cos\beta=2R(D+2R\cos\beta)=2R d_t.
\]
Введём отношение амплитуд нормальной и радиальной компонент
\begin{equation}
 \sigma=\operatorname{sgn}\Delta e_y,\qquad
 F=-\sigma\eta R,\qquad \eta>\sqrt3.
 \label{eq:ndc-eta-definition}
\end{equation}
Знак в этой записи соответствует рассматриваемому решению при
$\Delta i>0$; он будет проверен нормальными граничными условиями.
Поскольку $d_r=R\sin\beta$, условие касания теперь равно
\begin{equation}
 (1+\eta^2)d_r\cos\beta=2d_t\sin\beta,
 \qquad d_z=-\sigma\eta d_r.
 \label{eq:ndc-tangency}
\end{equation}
Ограничение $\eta>\sqrt3$ обеспечивает два максимума модуля
смещённого эллипса. Его смысл будет виден и при проверке годографа.

Введём сумму модулей и их относительную разность:
\begin{equation}
 \Delta\widehat V_+=\frac J2(1+\delta),\qquad
 \Delta\widehat V_-=\frac J2(1-\delta).
 \label{eq:ndc-split}
\end{equation}
Таким образом, $\delta$ определяет распределение затрат между
двумя импульсами. Неотрицательность обоих модулей равносильна
$|\delta|\leq1$.

Чтобы раскрыть синусы $\alpha\pm\beta$ без длинных выражений,
временно запишем компоненты изменения эксцентриситета в осях,
повёрнутых на $\alpha$:
\[
 \Delta e'_x=\Delta e_x\cos\alpha+\Delta e_y\sin\alpha,
 \qquad
 \Delta e'_y=-\Delta e_x\sin\alpha+\Delta e_y\cos\alpha.
\]
Штрих обозначает поворот осей, а не производную.
Подстановка~\eqref{eq:ndc-symmetry}--\eqref{eq:ndc-split}
в пять граничных уравнений даёт
\begin{equation}
\begin{aligned}
 \Delta a&=2Jd_t,\\
 \Delta e'_x&=J(d_r\sin\beta+2d_t\cos\beta),\\
 \Delta e'_y&=J\delta(2d_t\sin\beta-d_r\cos\beta),\\
 -\Delta i\sin\alpha&=Jd_z\sin\beta,\\
 \Delta i\cos\alpha&=J\delta d_z\cos\beta.
\end{aligned}
\label{eq:ndc-five-balances}
\end{equation}
Например, в первой строке компоненты $d_t$ одинаковы, поэтому
модули складываются. Во второй строке радиальная компонента
и синус одновременно меняют знак, поэтому они также дают сумму.
В третьей строке остаётся разность модулей $J\delta$.
Последние две строки --- повернутые условия~\eqref{eq:plane-change-boundaries}.

\subsubsection{Расчёт NDC: одно квадратное уравнение}

Исключение неизвестных из~\eqref{eq:ndc-tangency}
и~\eqref{eq:ndc-five-balances} приводит к следующему квадратному
уравнению для $\eta$:
\begin{equation}
 \Delta i|\Delta e_y|\eta^2-
 \bigl[(\Delta a)^2+(\Delta i)^2-\Delta e^2\bigr]\eta
 -\Delta i|\Delta e_y|=0.
 \label{eq:ndc-quadratic}
\end{equation}
Для удобства вычисления обозначим его средний коэффициент через
$\kappa=(\Delta a)^2+(\Delta i)^2-\Delta e^2$.
Положительный корень и дальнейшие параметры равны
\begin{equation}
\begin{aligned}
 \eta&=\frac{\kappa+\sqrt{\kappa^2+4(\Delta i)^2(\Delta e_y)^2}}
                  {2\Delta i|\Delta e_y|},\\
 L&=\sqrt{(\Delta e_x)^2+(|\Delta e_y|+\Delta i\eta)^2},\\
 \cos\alpha&=\frac{\Delta e_x}{L},\qquad
 \sin\alpha=\frac{\sigma(|\Delta e_y|+\Delta i\eta)}{L},\qquad
 \cos\beta=\frac{\Delta a}{L}.
\end{aligned}
\label{eq:ndc-angles}
\end{equation}
$L$ здесь --- длина двумерного вектора, определяющего $\alpha$;
она нужна только для компактной записи расчётных формул.
Угол $\alpha$ определяют через $\operatorname{atan2}$,
а $\beta=\arccos(\Delta a/L)$ выбирают в $(0,\pi)$.

\begin{explanation}{От системы~\eqref{eq:ndc-five-balances} к углам}
Из первой строки $Jd_t=\Delta a/2$.
Из условия касания
$Jd_r\cos\beta=\Delta a\sin\beta/(1+\eta^2)$.
Подставляя это и $d_z=-\sigma\eta d_r$ в третью и пятую строки,
получаем
\[
 \Delta e'_y=-\sigma\eta\Delta i\cos\alpha,
 \quad\text{то есть}\quad
 \Delta e_x\sin\alpha=(\Delta e_y+\sigma\eta\Delta i)\cos\alpha.
\]
Это даёт формулы для $\sin\alpha,\cos\alpha$ в~\eqref{eq:ndc-angles}.
Из второй и четвёртой строк аналогично следуют
\begin{align*}
 \Delta a L\cos\beta
    &=L^2-\Delta i(\eta+\eta^{-1})(|\Delta e_y|+\Delta i\eta),\\
 \Delta a L\frac{\sin^2\beta}{\cos\beta}
    &=\Delta i(\eta+\eta^{-1})(|\Delta e_y|+\Delta i\eta).
\end{align*}
Сложение даёт $\Delta a L/\cos\beta=L^2$, то есть
$\cos\beta=\Delta a/L$. Подстановка в первую строку даёт
\[
 (\Delta a)^2=L^2-\Delta i(\eta+\eta^{-1})(|\Delta e_y|+\Delta i\eta).
\]
Раскрывая $L^2$ и сокращая слагаемые, получаем
$\kappa=\Delta i|\Delta e_y|(\eta-\eta^{-1})$.
Умножение на $\eta$ даёт~\eqref{eq:ndc-quadratic}.
Деление на $\cos\beta$ допустимо во внутренности области NDC,
где $\Delta a\ne0$. Границы ниже отнесены к соседним семействам.
\end{explanation}

Теперь вычислим направления и суммарную стоимость.
Первые два соотношения в предыдущем блоке дают
\begin{equation}
 d_t=\frac{\Delta a}{2J},\qquad
 d_r=\frac{L\sin\beta}{(1+\eta^2)J},\qquad
 d_z=-\sigma\eta d_r.
 \label{eq:ndc-directions}
\end{equation}
Требование $d_r^2+d_t^2+d_z^2=1$ определяет $J$:
\begin{equation}
 \boxed{\ J_{\rm NDC}=
 \sqrt{\frac{(\Delta a)^2}{4}+
             \frac{L^2-(\Delta a)^2}{1+\eta^2}}.\ }
 \label{eq:ndc-cost}
\end{equation}
После исключения $L,\eta$ получается выражение Эдельбаума
\cite[с.~52, формула (68)]{edelbaum}:
\[
 J_{\rm NDC}^2=
 \frac{(\Delta i)^2+\Delta e^2-(\Delta a)^2/2+
       \sqrt{\kappa^2+4(\Delta i)^2(\Delta e_y)^2}}{2}.
\]
В практическом расчёте достаточно одной из этих двух равносильных формул.

Из последней строки~\eqref{eq:ndc-five-balances} находим разность
модулей, а затем сами импульсы:
\begin{equation}
\begin{aligned}
 \delta&=-\frac{\sigma\Delta e_x\sqrt{L^2-(\Delta a)^2}}
                    {\Delta a(|\Delta e_y|+\Delta i\eta)},\\
 \Delta\widehat{\boldsymbol V}_+
      &=\frac{J_{\rm NDC}}2(1+\delta)(d_r,d_t,d_z),\\
 \Delta\widehat{\boldsymbol V}_-
      &=\frac{J_{\rm NDC}}2(1-\delta)(-d_r,d_t,-d_z).
\end{aligned}
\label{eq:ndc-impulses}
\end{equation}
При восстановлении направлений в~\eqref{eq:ndc-directions}
берётся именно $J=J_{\rm NDC}$.

\paragraph{Область применимости и проверка минимума.}
Из квадратного уравнения следует
\[
 \eta>\sqrt3\quad\Longleftrightarrow\quad
 (\Delta a)^2>\Delta e^2+\frac{2\Delta i|\Delta e_y|}{\sqrt3}-(\Delta i)^2.
\]
Кроме того,
\[
 L^2-(\Delta a)^2=
 \frac{\Delta i(1+\eta^2)(|\Delta e_y|+\Delta i\eta)}{\eta}>0.
\]
Поэтому угол $\beta$ существует. Подстановка последнего тождества
в выражение для $\delta$ показывает, что $|\delta|\leq1$
равносильно $|\Delta a|\geq|\Delta e_x|$.
Во внутренности семейства принимаем строгие неравенства:
\begin{equation}
 \boxed{\begin{gathered}
 |\Delta a|>|\Delta e_x|,\\
 (\Delta a)^2>\Delta e^2+
            \frac{2\Delta i|\Delta e_y|}{\sqrt3}-(\Delta i)^2.
 \end{gathered}}
 \label{eq:ndc-region}
\end{equation}
Равенства на границах можно обрабатывать формулами NC или SC.

Для проверки годографа восстановим его коэффициенты:
\[
 R=\frac{L}{(1+\eta^2)J},\qquad F=-\sigma\eta R,
 \qquad D=\frac{\Delta a(\eta^2-3)}{2(1+\eta^2)J}.
\]
Вычисление квадрата модуля и выделение полного квадрата дают
\begin{equation}
 \|\boldsymbol s(\varphi)\|^2
 =1-(F^2-3R^2)\,[\cos(\varphi-\alpha)-\cos\beta]^2.
 \label{eq:ndc-norm}
\end{equation}
Так как $F^2-3R^2=R^2(\eta^2-3)>0$, годограф нигде не выходит
за сферу. Равенство достигается при $\varphi=\alpha\pm\beta$.
Граничные уравнения выполнены по построению, модули положительны,
направления совпадают с базис-вектором. Поэтому найденный переход
оптимален по критерию раздела~\ref{sec:optimality-check-full}
при $\lambda_6=0$.

\paragraph{Численный пример.}
При $\Delta a=0{,}004$, $\Delta e_x=\Delta e_y=0{,}001$,
$\Delta i=0{,}002$ получаем
\[
 J=0{,}00286697,\qquad
 \varphi_-=9{,}01739^\circ,\qquad\varphi_+=165{,}02574^\circ.
\]
Компоненты импульсов приведены в единицах $10^{-4}$:
\begin{center}
\begin{tabular}{c|rrr}
 Фаза, град & $\Delta\widehat V_r$ & $\Delta\widehat V_t$ & $\Delta\widehat V_z$\\\hline
 $9{,}01739$ & $-1{,}39514$ & $12{,}44871$ & $12{,}70944$\\
 $165{,}02574$ & $0{,}84628$ & $7{,}55129$ & $-7{,}70944$
\end{tabular}
\end{center}
Например, сумма трансверсальных компонентов равна $0{,}002$,
что обеспечивает $\Delta a=0{,}004$. При принятом $V_0$
размерная стоимость составляет $21{,}634$ м/с.

\subsubsection{Простой предел: годограф-отрезок}
\label{sec:line}

Если $\Delta e_y=0$, формула для $\eta$ содержит деление на нуль.
Сначала проверяют область NC. Если $|\Delta a|\leq|\Delta e_x|$,
применяется узловая схема. Если $|\Delta a|>|\Delta e_x|$,
можно взять годограф-отрезок
\begin{equation}
 J=\sqrt{\frac{(\Delta a)^2}{4}+(\Delta i)^2},\qquad
 \boldsymbol s(\varphi)=\frac1J(0,\Delta a/2,\Delta i\cos\varphi).
 \label{eq:line-primer}
\end{equation}
Его модуль не превосходит единицы, а концы отрезка при
$\varphi_1=0$, $\varphi_2=\pi$ лежат на сфере.
Трансверсальные направления теперь совпадают, а нормальные противоположны.

Подставим направления $\boldsymbol s(0),\boldsymbol s(\pi)$
в граничные уравнения. Условия на $\Delta a,\Delta i$ выполняются
при сумме модулей $J$; условие на $\Delta e_x$ задаёт разность.
В результате
\begin{equation}
\begin{aligned}
 \Delta\widehat V_1&=\frac J2\left(1+\frac{\Delta e_x}{\Delta a}\right),&
 \Delta\widehat V_2&=\frac J2\left(1-\frac{\Delta e_x}{\Delta a}\right),\\
 \Delta\widehat{\boldsymbol V}_1&=
       \frac{\Delta\widehat V_1}{J}(0,\Delta a/2,\Delta i),&
 \Delta\widehat{\boldsymbol V}_2&=
       \frac{\Delta\widehat V_2}{J}(0,\Delta a/2,-\Delta i).
\end{aligned}
\label{eq:line-impulses}
\end{equation}
Оба модуля положительны в указанной области. Радиальные условия
выполнены, поскольку $\Delta e_y=0$.
Особенно прост пример $\Delta e_x=\Delta e_y=0$:
два равных по модулю импульса одновременно меняют радиус круговой
орбиты и её плоскость. При $\Delta a=0$, $\Delta e_x=\Delta e_y=0$
достаточно одного нормального импульса $\Delta\widehat V_z=\Delta i$
в восходящем узле; это уже предельный случай NC.

\subsubsection{Особое решение: годограф целиком лежит на сфере}
\label{sec:singular}

В особом случае SC модуль базис-вектора равен единице при всех
фазах. В названии «особый» нет указания на ошибку или отсутствие
решения: исчезает однозначность выбора мест импульсов.
Направление в каждой фазе остаётся определённым.

При $\Delta i>0$ для этого семейства $\Delta e_y\ne0$.
Положим $\sigma=\operatorname{sgn}\Delta e_y$ и рассмотрим
\begin{equation}
 \boldsymbol s(\varphi)=
 \left(\tfrac12\sin\psi,\ \cos\psi,
          -\sigma\tfrac{\sqrt3}{2}\sin\psi\right),
 \qquad\psi=\varphi-\alpha.
 \label{eq:sc-primer}
\end{equation}
Квадраты радиальной и нормальной компонент в сумме равны
$\sin^2\psi$, поэтому $\|\boldsymbol s\|^2=\sin^2\psi+\cos^2\psi=1$.
Этот вектор имеет форму~\eqref{eq:primer-components}; соответствующие
множители равны
\[
 \lambda_1=0,\quad \lambda_3=\tfrac12\cos\alpha,\quad \lambda_2=\tfrac12\sin\alpha,
 \quad \lambda_4=-\sigma\tfrac{\sqrt3}{2}\cos\alpha,
 \quad \lambda_5=\sigma\tfrac{\sqrt3}{2}\sin\alpha.
\]
При $\lambda_6=0$ оценка~\eqref{eq:lower-bound-full} принимает вид
\[
 J\geq\frac12[\Delta e_x\cos\alpha+
                   (\Delta e_y+\sigma\sqrt3\Delta i)\sin\alpha].
\]
Выберем $(\cos\alpha,\sin\alpha)$ сонаправленным вектору
$(\Delta e_x,\Delta e_y+\sigma\sqrt3\Delta i)$.
Скалярное произведение тогда равно его длине, и получаем
\begin{equation}
\begin{aligned}
 J_s&=\frac12\sqrt{(\Delta e_x)^2+(|\Delta e_y|+\sqrt3\Delta i)^2},\\
 \cos\alpha&=\frac{\Delta e_x}{2J_s},\qquad
 \sin\alpha=\frac{\sigma(|\Delta e_y|+\sqrt3\Delta i)}{2J_s},
 \qquad J\geq J_s.
\end{aligned}
\label{eq:sc-cost}
\end{equation}
Пока $J_s$ является только нижней оценкой. Она станет достигаемым
минимумом, если удастся распределить импульсы суммарной величины $J_s$
по направлениям~\eqref{eq:sc-primer} так, чтобы выполнить все
граничные условия.

\paragraph{Какие уравнения остаётся решить.}
Подставим $\Delta\widehat{\boldsymbol V}_j
=\Delta\widehat V_j\boldsymbol s(\varphi_j)$ в граничные уравнения.
Как и в NDC, используем временно повернутые компоненты
$\Delta e'_x,\Delta e'_y$ с углом $\alpha$.
Получаются четыре простые суммы:
\begin{equation}
\begin{aligned}
 \sum_j\Delta\widehat V_j&=J_s,&
 \sum_j\Delta\widehat V_j\cos\psi_j&=\frac{\Delta a}{2},\\
 \sum_j\Delta\widehat V_j\sin^2\psi_j&=\Sigma_{ss},&
 \sum_j\Delta\widehat V_j\sin\psi_j\cos\psi_j&=\Sigma_{sc},
\end{aligned}
\label{eq:sc-four-balances}
\end{equation}
где две вспомогательные суммы заранее известны:
\begin{equation}
 \Sigma_{ss}=\frac{\Delta i(|\Delta e_y|+\sqrt3\Delta i)}{\sqrt3J_s},
 \qquad \Sigma_{sc}=-\frac{\sigma\Delta i\Delta e_x}{\sqrt3J_s}.
 \label{eq:sc-uv}
\end{equation}
Индексы у $\Sigma_{ss}$ и $\Sigma_{sc}$ указывают на множители
$\sin\psi\sin\psi$ и $\sin\psi\cos\psi$ в соответствующих суммах.

\begin{explanation}{Почему четырёх сумм~\eqref{eq:sc-four-balances} достаточно}
После поворота осей изменения эксцентриситета равны
\[
 \Delta e'_x=\sum_j\Delta\widehat V_j
       (\tfrac12\sin^2\psi_j+2\cos^2\psi_j)
       =2J_s-\tfrac32\Sigma_{ss},
 \qquad \Delta e'_y=\tfrac32\Sigma_{sc}.
\]
Нормальные условия в этих же осях имеют вид
\[
 -\sigma\tfrac{\sqrt3}{2}\Sigma_{ss}=-\Delta i\sin\alpha,
 \qquad -\sigma\tfrac{\sqrt3}{2}\Sigma_{sc}=\Delta i\cos\alpha.
\]
Последние две формулы дают~\eqref{eq:sc-uv}.
Подстановка $\alpha,J_s$ из~\eqref{eq:sc-cost} показывает, что они
одновременно удовлетворяют обоим условиям на эксцентриситет.
Условие на большую полуось --- вторая сумма
в~\eqref{eq:sc-four-balances}. Таким образом, выполнены все пять
орбитальных условий и назначена суммарная стоимость.
\end{explanation}

\paragraph{Выбор мест по построению Эдельбаума.}
Следуя~\cite[с.~56]{edelbaum}, выберем одинаковое значение
$\sin^2\psi_j$ для всех импульсов. Обозначим
\begin{equation}
 s=\sqrt{\Sigma_{ss}/J_s},\qquad c=\sqrt{1-s^2},\qquad
 \beta=\operatorname{atan2}(s,c).
 \label{eq:sc-angle}
\end{equation}
Эти $s,c$ --- сокращённая запись синуса и косинуса вспомогательного
угла $\beta$ только в данном расчёте.
При $0<s,c<1$ допустимы четыре положения:
\begin{equation}
 \psi_1=\pi-\beta,\qquad \psi_2=-\beta,\qquad
 \psi_3=\beta,\qquad \psi_4=\pi+\beta.
 \label{eq:sc-four-angles}
\end{equation}
Условие с $\sin^2\psi_j$ теперь выполнено автоматически:
его левая часть равна $s^2\sum_j\Delta\widehat V_j=s^2J_s=\Sigma_{ss}$.

\begin{figure}[htbp]
\centering
\begin{tikzpicture}[scale=2,>=Latex]
 \draw[->,gray] (-1.25,0)--(1.45,0) node[right] {$\cos\psi$};
 \draw[->,gray] (0,-1.25)--(0,1.3) node[above] {$\sin\psi$};
 \draw[accent,thick] (0,0) circle(1);
 \draw[dashed] (-1.1,.6)--(1.1,.6) node[right] {$s$};
 \draw[dashed] (-1.1,-.6)--(1.1,-.6) node[right] {$-s$};
 \foreach \xx/\yy/\name in {-.8/.6/{1},.8/-.6/{2},.8/.6/{3},-.8/-.6/{4}}{
  \fill[accent] (\xx,\yy) circle(1.4pt);
  \draw[accent] (0,0)--(\xx,\yy);}
 \node[above left] at (-.8,.6) {$\psi_1$};
 \node[below right] at (.8,-.6) {$\psi_2$};
 \node[above right] at (.8,.6) {$\psi_3$};
 \node[below left] at (-.8,-.6) {$\psi_4$};
 \draw[->] (.4,0) arc[start angle=0,end angle=36.87,radius=.4];
 \node at (.5,.16) {$\beta$};
\end{tikzpicture}
\caption{Четыре вспомогательных угла с одинаковым квадратом синуса.
Физические места импульсов имеют фазы $\varphi_j=\alpha+\psi_j$.
Для решения достаточно выбрать не более трёх из четырёх мест.}
\label{fig:sc-four-angles}
\end{figure}

Остаются сумма модулей и два уравнения, зависящих только от знаков
синуса и косинуса. Их удобно прочитать из таблицы:
\begin{center}
\begin{tabular}{c|rrrr}
 Номер места &1&2&3&4\\\hline
 $\cos\psi_j$ & $-c$ & $c$ & $c$ & $-c$\\
 $\sin\psi_j\cos\psi_j$ & $-sc$ & $-sc$ & $sc$ & $sc$
\end{tabular}
\end{center}
Поэтому система для четырёх модулей равна
\begin{equation}
\begin{aligned}
 \Delta\widehat V_1+\Delta\widehat V_2+\Delta\widehat V_3+\Delta\widehat V_4&=J_s,\\
 -\Delta\widehat V_1+\Delta\widehat V_2+\Delta\widehat V_3-\Delta\widehat V_4
       &=\frac{\Delta a}{2c},\\
 -\Delta\widehat V_1-\Delta\widehat V_2+\Delta\widehat V_3+\Delta\widehat V_4
       &=\frac{\Sigma_{sc}}{sc}.
\end{aligned}
\label{eq:sc-linear-system}
\end{equation}
Это три линейных уравнения с четырьмя неизвестными.
Одно неизвестное можно обнулить, выбрав вариант, при котором
остальные три неотрицательны.

Если $2\Sigma_{sc}\leq\Delta a\,s$, полагаем $\Delta\widehat V_4=0$:
\begin{equation}
 \Delta\widehat V_1=\frac{J_s}{2}-\frac{\Delta a}{4c},\qquad
 \Delta\widehat V_2=\frac{\Delta a}{4c}-\frac{\Sigma_{sc}}{2sc},\qquad
 \Delta\widehat V_3=\frac{J_s}{2}+\frac{\Sigma_{sc}}{2sc}.
 \label{eq:sc-first-three}
\end{equation}
Если $2\Sigma_{sc}\geq\Delta a\,s$, полагаем $\Delta\widehat V_2=0$:
\begin{equation}
 \Delta\widehat V_1=\frac{J_s}{2}-\frac{\Sigma_{sc}}{2sc},\qquad
 \Delta\widehat V_3=\frac{J_s}{2}+\frac{\Delta a}{4c},\qquad
 \Delta\widehat V_4=\frac{\Sigma_{sc}}{2sc}-\frac{\Delta a}{4c}.
 \label{eq:sc-last-three}
\end{equation}
При равенстве оба варианта совпадают и остаются не более двух
ненулевых импульсов. Окончательные векторы:
\begin{equation}
 \varphi_j=\alpha+\psi_j\pmod{2\pi},\qquad
 \Delta\widehat{\boldsymbol V}_j=
       \Delta\widehat V_j\boldsymbol s(\varphi_j).
 \label{eq:sc-final-impulses}
\end{equation}
Нулевые импульсы исключают, остальные сортируют по фазе.

\paragraph{Когда построение допустимо.}
Из~\eqref{eq:sc-linear-system} видно, что при неотрицательных модулях
\[
 |\Delta a|\leq2J_sc,\qquad |\Sigma_{sc}|\leq J_ssc.
\]
Эти же условия обеспечивают неотрицательность всех модулей
в выбранной тройке~\eqref{eq:sc-first-three} или~\eqref{eq:sc-last-three}.
Выразим ограничения через исходные орбиты:
\begin{equation}
 \boxed{\begin{gathered}
 \Delta i\leq\sqrt3\,|\Delta e_y|,\\
 (\Delta a)^2\leq\Delta e^2+
         \frac{2\Delta i|\Delta e_y|}{\sqrt3}-(\Delta i)^2.
 \end{gathered}}
 \label{eq:sc-region}
\end{equation}

\begin{explanation}{Переход к области~\eqref{eq:sc-region}}
Из определения $c$ следует
\[
 4J_s^2c^2=\Delta e^2+
       \frac{2\Delta i|\Delta e_y|}{\sqrt3}-(\Delta i)^2.
\]
Поэтому $|\Delta a|\leq2J_sc$ даёт вторую строку области.
Для второго ограничения вычислим
\[
 s^2c^2-\frac{\Sigma_{sc}^2}{J_s^2}
   =\frac{\Delta i(\sqrt3|\Delta e_y|-\Delta i)}{3J_s^2}.
\]
Левая часть должна быть неотрицательной. При $\Delta i>0$
это даёт первую строку области. Внутри неё $s,c$ действительны,
а формулы трёх модулей удовлетворяют всем ограничениям.
\end{explanation}

Построенные импульсы достигают нижней оценки $J_s$, значит,
переход оптимален. Отсутствие $\Delta a$ в формуле стоимости
не означает, что большая полуось не меняется: её изменение
обеспечивает перераспределение модулей между выбранными местами.
Это возможно только в области~\eqref{eq:sc-region}.

Построение использует не более трёх импульсов. Эдельбаум также
показывает существование двухимпульсных оптимальных реализаций
особого случая~\cite[с.~55--56]{edelbaum}, но их явный расчёт сложнее.
Трёхимпульсная схема имеет ту же минимальную характеристическую
скорость. Здесь минимизируется сумма модулей, а не число включений.

При $c=0$ деление в~\eqref{eq:sc-linear-system} недопустимо.
Из условий области тогда следуют
$\Delta a=\Delta e_x=0$, $\Delta i=\sqrt3|\Delta e_y|$.
Это общая граница с NC, где достаточно одного импульса
$(-\Delta e_y,0,\Delta i)$ в узле. При $s=0$ имеем
$\Delta i=0$ и применяем компланарный алгоритм.

\paragraph{Численный пример.}
Пусть $\Delta a=0{,}001$, $\Delta e_x=0$,
$\Delta e_y=0{,}001$, $\Delta i=0{,}0002$.
Тогда $J_s=0{,}000673205$, $\alpha=90^\circ$,
$\beta=35{,}85253^\circ$ и $\Sigma_{sc}=0$.
Схема~\eqref{eq:sc-first-three} после сортировки фаз даёт
следующие компоненты в единицах $10^{-4}$:
\begin{center}
\begin{tabular}{c|rrr}
 Фаза, град & $\Delta\widehat V_r$ & $\Delta\widehat V_t$ & $\Delta\widehat V_z$\\\hline
 $54{,}14747$ & $-0{,}90327$ & $2{,}50000$ & $1{,}56451$\\
 $125{,}85253$ & $0{,}98574$ & $2{,}72826$ & $-1{,}70736$\\
 $234{,}14747$ & $0{,}08247$ & $-0{,}22826$ & $-0{,}14284$
\end{tabular}
\end{center}
Сумма трансверсальных компонентов равна $0{,}0005$, как требуется
для $\Delta a=0{,}001$. При принятом $V_0$ стоимость составляет
$5{,}080$ м/с.

\subsection{Порядок расчёта и проверка результата}
\label{sec:transfer-check}

Классификацию удобно применять в следующем порядке. Это одновременно
устраняет деление на нуль в формулах отдельных семейств.
\begin{enumerate}
 \item Если $\Delta i=0$, использовать компланарный алгоритм.
 При нулевых $\Delta a,\Delta e$ импульсы отсутствуют.
 \item Если $|\Delta a|\leq|\Delta e_x|$ и
 $\Delta i\geq\sqrt3|\Delta e_y|$, использовать NC.
 При $\Delta e_x=0$ применяется отдельная одноимпульсная формула.
 \item Если $\Delta e_y=0$ и NC не подходит, использовать
 годограф-отрезок~\eqref{eq:line-primer}--\eqref{eq:line-impulses}.
 \item Если выполнены оба неравенства~\eqref{eq:sc-region},
 использовать построение SC.
 \item В остальных случаях использовать NDC:
 решить~\eqref{eq:ndc-quadratic}, найти углы, направления и модули
 по~\eqref{eq:ndc-angles}--\eqref{eq:ndc-impulses}.
\end{enumerate}

Почему случаев достаточно? Разность двух границ по $(\Delta a)^2$
разлагается на множители:
\[
 \Delta e^2+\frac{2\Delta i|\Delta e_y|}{\sqrt3}-(\Delta i)^2
       -(\Delta e_x)^2
 =\left(|\Delta e_y|-\frac{\Delta i}{\sqrt3}\right)
       \left(|\Delta e_y|+\sqrt3\Delta i\right).
\]
При $\Delta i\geq\sqrt3|\Delta e_y|$ первой достигается граница NC,
а при обратном неравенстве --- граница SC. Если соответствующая
граница по $|\Delta a|$ превышена, выполнены условия NDC.
На общих границах стоимости совпадают, а оптимальная схема
может иметь несколько представлений.

\begin{explanation}{Что проверять после расчёта}
По найденным фазам и компонентам независимо вычисляют левые части
всех пяти уравнений~\eqref{eq:opt-boundaries} и сравнивают с
$\Delta a,\Delta e_x,\Delta e_y,0,\Delta i$.
Далее независимо вычисляют сумму модулей:
\[
 J=\sum_j\sqrt{(\Delta\widehat V_{rj})^2+
(\Delta\widehat V_{tj})^2+(\Delta\widehat V_{zj})^2}.
\]
Её сравнивают с формулой выбранного семейства.
Отрицательной может быть компонента; отрицательный модуль
означает, что выбранное семейство или распределение неприменимо.
Проверка ограничений на годограф уже доказана аналитически
для каждой указанной области.
\end{explanation}

Фазы сначала определяются по модулю $2\pi$. Для выполнения манёвра
выбирают их хронологический порядок: после импульса в фазе
$\varphi_1$ следующий должен иметь больший развёрнутый угол.
При необходимости к фазе добавляют $2\pi$. В принятой модели
пассивный интервал равен $\Delta t=(\varphi_2-\varphi_1)/n_0$.
Перестановка манёвров может менять промежуточную орбиту и длительность,
но сохраняет пять суммарных условий перехода и сумму модулей.

Приведённые алгоритмы решают задачу с любыми разрешёнными фазами,
свободной длительностью и без ограничений на промежуточную орбиту.
Для встречи в заданный момент нужно дополнительно выполнить фазовое
условие из главы~1; его нельзя автоматически считать выполненным
после решения пяти уравнений перехода.
